% Printed 401
\par\clearpage
\noindent\null\hfill{\small\textbf{ZUR HANNOVERSCHEN TRIANGULATION.}}\hfill\makebox[0pt][r]{\fontsize{13}{16}\selectfont 401}\par
\vspace*{4\baselineskip}
\begin{center}
{\large III.}\par\medskip
{\LARGE\textls[150]{NACHLASS}}\par\smallskip
{\large ZUR HANNOVERSCHEN TRIANGULATION.}
\end{center}
\vspace{2\baselineskip}
\begin{center}
[1.]\par\smallskip
{\large [Plan und Anfang des Werks über die trigonometrischen Messungen\\
in Hannover.]}\par\medskip
\textls[120]{Plan des Werks.}\par\smallskip
Acht Abschnitte.
\end{center}

\noindent
\begin{tabular*}{\linewidth}{@{}r@{\enspace}l@{\extracolsep{\fill}}r@{\hspace{3.6em}}}
1) & Über die bei den Messungen angewandten Einrichtungen & \\
   & und Methoden im Allgemeinen & \makebox[0pt][r]{ungefähr\enspace}40\makebox[0pt][l]{\enspace Seiten}\\
2) & Messung der Hauptdreieckswinkel & 60\\
3) & Kleine Reduktionen der Winkelmessungen & 12\\
4) & Ausgleichung der Messungen & 24\\
5) & Erste Berechnungsart der Dreiecke, in der Darstellung & \\
   & auf dem Sphäroid & 36\\
6) & Zweite Berechnungsart, in der Darstellung in der Ebene & 60\\
7) & Nebenpunkte & 24\\
8) & Höhenmessungen & 20\\
\multicolumn{3}{@{}r@{}}{Zusammen etwa 276 Seiten}\\
\multicolumn{3}{@{}r@{}}{oder \GaussIXMathCopy{34 1/2}{\(34\tfrac{1}{2}\)} Bogen.}
\end{tabular*}\par\medskip

\noindent\hspace*{\parindent}Im ersten Abschnitt sind folgende Gegenstände zu behandeln: unmittelbarer
Zweck der Messungen; Eigenart des Geländes; Zielpunkte; Standpunkte;
Instrumente; Beobachtungsmethode; Basis; vorläufige Berechnung der Punkte.
\par\smallskip
\noindent\hspace*{.031\linewidth}{\footnotesize IX.}\hfill{\small 51}\hspace*{.128\linewidth}\par

% Printed 402
\clearpage
\noindent\makebox[0pt][l]{\fontsize{13}{16}\selectfont 402}\hfill{\small\textbf{NACHLASS.}}\hfill\null\par
\vspace{2.5\baselineskip}
\begin{center}
{\large\textls[120]{Die trigonometrischen Messungen im Königreich Hannover.}}\par\medskip
\rule{.12\linewidth}{.4pt}\par\medskip
{\large\textls[120]{Einleitung.}}
\end{center}

\par\noindent\hspace*{\parindent}Den ersten Anlass zu den trigonometrischen Messungen, die ich in den Jahren
1821---1825 im Königreich Hannover ausführte, gab die von Herrn Etatsrat
\textsc{Schumacher} in den dänischen Staaten unternommene Gradmessung. Während
diese an der Nordspitze Jütlands ihren natürlichen Endpunkt finden musste, ließ
sie sich nach Süden ausdehnen und würde erst am Mittelmeer an ihre Grenze stoßen. Eine solche Erweiterung musste
zunächst über eine Strecke von mehr als zwei Breitengraden durch das Königreich
Hannover führen. An dessen Südgrenze bot die Göttinger Sternwarte Gelegenheit
zu den genauesten und bequemsten astronomischen Bestimmungen. Die
wissenschaftliche Bedeutung eines solchen Unternehmens und die feste Grundlage,
die dadurch die Geografie des Königreichs und künftige umfassendere
trigonometrische Messungen erhalten mussten, konnten unserer aufgeklärten
Regierung nicht entgehen. Daher erhielt ich 1820 den Auftrag, diese Fortsetzung
der dänischen Gradmessung durch das Königreich Hannover auszuführen.

\par\noindent\hspace*{\parindent}Der trigonometrische Teil dieser Arbeit wurde im Sommer 1821 am
südlichen Ende begonnen und im Sommer 1823 mit den Messungen auf dem
Michaelisthurm in Hamburg abgeschlossen.

\par\noindent\hspace*{\parindent}Inzwischen hatte im benachbarten Kurfürstentum Hessen unter der Leitung
von Herrn Professor \textsc{Gerling} eine trigonometrische Landesvermessung begonnen,
bei der die Hauptdreiecke mit ausgezeichneten Hilfsmitteln und mit aller
erreichbaren Genauigkeit gemessen werden sollten. Ihre Verbindung mit den
hannoverschen Dreiecken war daher umso wichtiger, weil diese dadurch mit den
bayerischen Dreiecken in Zusammenhang kommen mussten und die in verschiedenen
Teilen Europas ausgeführten Dreiecksmessungen durch ihre Verknüpfung zu einem
einzigen Ganzen in einem übergeordneten
\par\clearpage

% Printed 403
\noindent\null\hfill{\small\textbf{ZUR HANNOVERSCHEN TRIANGULATION.}}\hfill\makebox[0pt][r]{\fontsize{13}{16}\selectfont 403}\par\medskip
\noindent wissenschaftlichen Zusammenhang ein Vielfaches ihres Werts erhalten würden.
Die Verbindung der hannoverschen und kurhessischen Messungen wurde noch
im Spätjahr 1823 ausgeführt. Letztere sind aber seitdem unvollendet geblieben,
obwohl das für eine notdürftige Verbindung der hannoverschen mit den bayerischen
Dreiecken Erforderliche, nämlich die Messung von wenigstens zwei Winkeln
in allen für die Verbindung nötigen Dreiecken, 1823 abgeschlossen wurde.

\par\noindent\hspace*{\parindent}Nachdem so eine wenigstens vorläufige Verknüpfung der nord- und
süddeutschen Messungen erreicht war, musste es doppelt wichtig erscheinen,
auch eine Verknüpfung mit den großen zusammenhängenden Messsystemen im
Westen herzustellen. Mein ursprünglicher Auftrag wurde deshalb dahingehend
erweitert, dass ich auch noch einen Anschluss meiner Gradmessungsdreiecke an
die Messungen \textsc{Krayenhoff}s herstellen sollte. Dazu legte ich in den Jahren
1824 und 1825 ein neues Dreieckssystem von Hamburg bis Jever an und stellte
so den Zusammenhang mit den niederländischen, französischen und englischen
Dreiecken her. Damit liegen tatsächlich schon jetzt alle großen Messungen vor,
die sich über den kultiviertesten Teil ganz Europas erstrecken.

\par\noindent\hspace*{\parindent}Da die dänische Gradmessung inzwischen noch unvollendet geblieben
war, hielt ich es für notwendig, den astronomischen Teil meiner Arbeit so
einzurichten, dass die hannoversche Gradmessung auch als ein abgeschlossenes
Ganzes für sich bestehen konnte.

\par\noindent\hspace*{\parindent}Die Göttinger Sternwarte, die selbst ein Hauptdreieckspunkt im System
ist und von der aus alle Dreiecksseiten ihre Orientierung erhielten, bildete von
selbst den südlichen Endpunkt. Ein in seiner Art einzigartiger Umstand ließ
auch an der Wahl des nördlichen Endpunkts keinen Zweifel: Die inzwischen in
Altona errichtete Sternwarte von Herrn Professor \textsc{Schumacher} liegt nämlich
fast genau auf dem Meridian der Göttinger Sternwarte. Die Beobachtungen,
die ich 1827 an beiden Orten mit dem \textsc{Ramsden}schen Zenitsektor an
43 Sternen angestellt habe, sowie den daraus erhaltenen Breitenunterschied und die
übrigen Ergebnisse habe ich bereits in einem 1828 erschienenen Werk veröffentlicht.

\par\noindent\hspace*{\parindent}Dem trigonometrischen Teil meiner Arbeit ist das vorliegende Werk ge-
\par\smallskip
\noindent\hfill{\small 51\textsuperscript{*}}\hspace*{.122\linewidth}\par\clearpage

% Printed 404
\noindent\makebox[0pt][l]{\fontsize{13}{16}\selectfont 404}\hfill{\small\textbf{NACHLASS.}}\hfill\null\par\medskip
\noindent widmet. Aus zwei Gründen hielt ich es dabei für nötig, der Darstellung
alle erforderliche Ausführlichkeit zu geben.

\par\noindent\hspace*{\parindent}Bei einer \textls[120]{isolierten} Breitengradmessung von mäßiger Ausdehnung,
\textls[120]{die nichts weiter als eine solche ist und sein soll}, kann man den
trigonometrischen Teil gewissermaßen als untergeordnet betrachten. Die
erreichbare Genauigkeit der astronomischen Beobachtungen entspricht bei
Weitem nicht der beim trigonometrischen Teil erreichbaren Genauigkeit.
Es ist also weniger unerlässlich, dabei die höchstmögliche Genauigkeit zu erreichen.

\par\noindent\hspace*{\parindent}Bei der trigonometrischen Vermessung eines Landes ist es dagegen in
mehrfacher Hinsicht ratsam, die gegenseitige Lage der Hauptpunkte so genau
zu bestimmen, wie es der Stand der Technik und die Umstände zulassen.
Dann kann man nämlich in unzähligen Fällen hinreichend genaue abgeleitete
Bestimmungen sekundärer Punkte mit äußerst geringem Aufwand und mit Methoden
gewinnen, die ohne jene Voraussetzung zu völlig unbrauchbaren Ergebnissen führen
würden. Steht eine solche trigonometrische Vermessung für sich, ist die
ausführlichere Veröffentlichung ihrer Bestandteile freilich nicht von allgemeinem
Interesse. Doch je mehr die in verschiedenen Teilen Europas ausgeführten
trigonometrischen Messungen miteinander verbunden werden und sich nach und nach
einem großen Ganzen nähern, desto mehr gewinnen die einzelnen Bestandteile
den Charakter eines kostbaren Gemeinguts von bleibendem Wert für alle Zeiten,
und desto wichtiger wird es, alle ihre wesentlichen Einzelheiten so vollständig
zu bewahren, dass ihre Zuverlässigkeit im Ganzen wie im Einzelnen stets
geprüft werden kann.

\par\noindent\hspace*{\parindent}Ein zweiter Grund für die Ausführlichkeit lag in der Eigenart der
Methoden, die ich sowohl bei den Messungen selbst als auch bei ihrer Verarbeitung
zu Ergebnissen anwandte. Sie unterscheiden sich zum Teil völlig von den
sonst üblichen Methoden, und ihre Darstellung sollte ein Hauptzweck dieses
Werks sein. Die allgemeine Darstellung dieser Methode mit ihrer fortlaufenden
Anwendung auf ein ausgedehntes Messsystem zu verbinden, ist zweifellos das
geeignetste Mittel, ihre Natur im richtigen Licht zu zeigen und denen die
Anwendung zu erleichtern, die sich ihrer künftig bei ähnlichen Messungen
bedienen wollen.
\par\clearpage

% Printed 405
\noindent\null\hfill{\small\textbf{ZUR HANNOVERSCHEN TRIANGULATION.}}\hfill\makebox[0pt][r]{\fontsize{13}{16}\selectfont 405}\par\medskip
\begin{center}
\textls[120]{Erster Abschnitt.}\par\smallskip
Anordnung der Messungen im Allgemeinen.\par\medskip
1.
\end{center}

\par\noindent\hspace*{\parindent}Der Landstrich von Göttingen bis Hamburg ist im Süden und im Norden
von sehr unterschiedlicher Beschaffenheit. Der Süden ist gebirgig. Die
Berggipfel sind meist mehr oder weniger bewaldet, und die meisten Ortschaften
liegen so, dass ihre Türme entweder gar keine weite Aussicht bieten oder
höchstens in eine Richtung. Der Norden dagegen ist flach und vielfach von
Wald durchzogen, was die Nutzung einzelner niedriger Anhöhen sehr erschwert
und oft ganz unmöglich macht. Ortschaften mit Türmen, die sich als
Dreieckspunkte eignen, fehlen auf dem größten Teil dieser Strecke völlig.

\begin{center}
2.
\end{center}

\par\noindent\hspace*{\parindent}Unter diesen Umständen ließ sich voraussehen, dass man kaum mit
natürlichen Dreieckspunkten rechnen konnte und an den meisten Dreieckspunkten
entweder eigens Signaltürme errichten oder andere künstliche Mittel in Betracht
ziehen musste, um sie sichtbar zu machen.

\par\noindent\hspace*{\parindent}Eine wesentliche Frage ist dabei, wie groß die einzelnen Dreiecke sein
sollen. Bei einer sehr groß angelegten Messung, die z.~B. einen ganzen Kontinent
umfasst, wäre es für die Genauigkeit des Ganzen im Allgemeinen am günstigsten,
die Dreiecke so groß wie möglich zu machen. Dasselbe gilt für eine Messung
von kleinerem Umfang, sofern man sie als Bestandteil eines solchen größeren
Systems betrachtet. Diese Behauptung trifft jedoch nur unter der Voraussetzung
zu, dass die Winkel in den größten Dreiecken mit derselben oder zumindest
mit keiner erheblich geringeren Präzision gemessen werden können als in
kleinen Dreiecken. Diese Bedingung ist bei der Verwendung von Kirchtürmen
oder künstlichen Signaltürmen keineswegs erfüllt. Augenblicke, in denen solche
Objekte in sehr großen Entfernungen die größte Präzision und Sicherheit
bei den Messungen erlauben, sind äußerst selten.

\begin{center}
\rule{.24\linewidth}{.4pt}
\end{center}
\clearpage

% Printed 406
\noindent\makebox[0pt][l]{\fontsize{13}{16}\selectfont 406}\hfill{\small\textbf{NACHLASS.}}\hfill\null\par
\vspace{2\baselineskip}
\begin{center}
[2.]\par\smallskip
{\large [Auszüge aus Berichten über die Triangulation an das hannoversche\\
Kabinettsministerium.]}
\end{center}
\bigskip
\begin{center}
[Aus einem Bericht vom 7.~Januar 1822 »über die Arbeiten im Jahre 1821«.]
\end{center}

\par\noindent\hspace*{\parindent}Von den zahlreichen Arbeiten, die zu einer Gradmessung gehören,
nimmt die Anlage des Dreiecksnetzes zusammen mit der Messung der Winkel
bei Weitem die meiste Zeit und Arbeit in Anspruch. Bei der hannoverschen Gradmessung muss
dieses Dreiecksnetz im Norden bei Hamburg an die dänischen Messungen
anschließen. Im Süden ist zwar die Göttinger Sternwarte der eigentliche
natürliche Endpunkt der Gradmessung an sich. Damit sie sich aber weiter nach
Süden ausdehnen lässt, ist es zugleich sehr wichtig, sie an diejenigen auswärtigen
Messungen anzuschließen, die das Königreich Hannover im Süden berühren.
Dazu gehören die von der königlichen preußischen Regierung veranlassten und
mit großer Sorgfalt ausgeführten Messungen. Auch in Kurhessen ist derzeit
eine große Triangulation geplant, die mit aller erreichbaren Genauigkeit
ausgeführt werden soll.

\par\noindent\hspace*{\parindent}Die erwähnten preußischen Messungen wurden mir im vergangenen Winter
mitgeteilt, soweit sie bis dahin gediehen waren. Außerdem hatte ich Gelegenheit,
einen Teil der Messungen zu erhalten, die der französische Oberst \textsc{Epailly}
im Jahr 1804 und in den folgenden Jahren im Hannoverschen, insbesondere in
dessen südlichem Teil, vorgenommen hatte. Diese und einige weitere Hilfsmittel
konnten mir die erste Auswahl der Dreieckspunkte erleichtern. Außerdem
bedachte ich, dass sich manche kleineren Schwierigkeiten und Probleme, die
vom Beginn solcher Arbeiten nicht zu trennen sind, in der Nähe von Göttingen
immer schneller und leichter beheben ließen. Das bewog mich, mit der
Triangulation \textls[120]{auf der Südseite} zu beginnen.

\par\noindent\hspace*{\parindent}Schon 1820 hatte ich geeignete Vorbereitungen getroffen, um mir die
erforderlichen Instrumente zu beschaffen\textperiodcentered{} hier erwähne ich nur der-
\par\clearpage

% Printed 407
\noindent\null\hfill{\small\textbf{ZUR HANNOVERSCHEN TRIANGULATION.}}\hfill\makebox[0pt][r]{\fontsize{13}{16}\selectfont 407}\par\medskip
\noindent jenigen, die unmittelbar den geodätischen Teil der Gradmessung betreffen.
Für die eigentlichen Winkelmessungen hatte ich bei \textsc{Reichenbach} (dessen
Werkstatt sein ehemaliger Werkmeister \textsc{Ertel} mittlerweile vollständig
übernommen hat) einen zwölfzölligen Theodoliten bestellt, dessen Fertigstellung
und Lieferung für das Frühjahr 1821 zugesagt waren. Einen kleineren
Theodoliten des englischen Instrumentenbauers \textsc{Troughton} hatte mir Professor
\textsc{Schumacher} freundlicherweise bereits besorgt.

\par\noindent\hspace*{\parindent}Die Hilfsmittel, um die Dreieckspunkte aus sehr großer Entfernung
sichtbar zu machen, erforderten besondere Vorkehrungen. Da ich die Dreiecke
so groß wie möglich wählen wollte und dies von größter Bedeutung war, ließ
die Beschaffenheit der Gegend, durch die sie verlaufen sollten, keine Hoffnung,
viele Kirchtürme als Dreieckspunkte verwenden zu können. Eigens errichtete
Signaltürme waren bisher das in solchen Fällen am häufigsten verwendete
Mittel. In der Praxis kommt es jedoch nicht selten vor, dass auch dieses Mittel
nicht ausreicht. Solche Signaltürme sind ebenso wie die Kirchtürme aus großer
Entfernung immer sehr schwer zu sehen und noch viel schwerer zu beobachten,
wenn sie sich nicht gegen den Himmel abheben, besonders wenn sie vor einem
nahen dunkelfarbigen Hintergrund erscheinen\textsuperscript{*}). Aus diesen und anderen
Gründen haben andere Beobachter die Winkel häufig bei Nacht beobachtet
(manche ausschließlich) und die entfernten Dreieckspunkte durch große
\textsc{Argand}sche Lampen mit sehr genau parabolischen Reflektoren sichtbar
machen lassen. Freilich bringen diese nächtlichen Beobachtungen wiederum
andere große Schwierigkeiten und Unannehmlichkeiten mit sich. Besonders
bei sehr großen Dreiecken muss eine gelungene Beobachtung gewöhnlich
gleichsam mit vielen vergeblichen Versuchen erkauft \mbox{werden.\enspace\mbox{.\ .\ .\ .}}\quad Obwohl
ich daher nicht dazu neigte, ausschließlich auf diese Art zu beobachten, zumal
meine körperlichen Kräfte den Beschwerden ständiger nächtlicher Aufenthalte
auf meist hohen und schwer zugänglichen Bergen wohl kaum gewachsen
gewesen wären, musste ich mich doch,

\par\smallskip\noindent\rule{.24\linewidth}{.4pt}\par\smallskip
{\footnotesize
\noindent\hspace*{\parindent}\textsuperscript{*}) Meine eigenen Erfahrungen im vergangenen Sommer haben dies vielfach bestätigt.
Während meines gesamten, mehr als vierwöchigen Aufenthalts auf dem Brocken habe ich z.~B. den auf dem Hils
errichteten, \(7\tfrac{1}{2}\) Meilen entfernten Signalturm nur ein- oder zweimal für wenige Minuten sehen können.
Die Kirchtürme des 12 Meilen entfernten Hannover konnte ich kein einziges Mal sehen, obwohl die Richtung genau bekannt war.\par}
\clearpage

% Printed 408
\noindent\makebox[0pt][l]{\fontsize{13}{16}\selectfont 408}\hfill{\small\textbf{NACHLASS.}}\hfill\null\par\medskip
\noindent da bisher kein anderes Mittel bekannt war, von vornherein darauf
einstellen, es wenigstens in manchen Einzelfällen zu benutzen. Daher hatte ich
vorläufig drei solcher Lampen bei \textsc{Repsold} in Hamburg und bei \textsc{Körner} in
Jena bestellt. Ich erhielt diese Lampen im Mai 1821, und ihre Wirkung bei
den verschiedenen Versuchen, die ich damit in geeigneten Entfernungen
anstellte, entsprach auch meiner Erwartung.

\par\noindent\hspace*{\parindent}Die Erfahrungen anderer führten mir schon vor meinen eigenen
Erfahrungen all die genannten großen Schwierigkeiten bei der Bildung großer
Dreiecke vor Augen, und ich suchte nach einem ganz neuen Mittel, um sie zu
überwinden. Theoretische Untersuchungen hatten mich davon überzeugt, dass
das selbst von ganz kleinen Planspiegeln reflektierte Sonnenlicht stark genug ist,
um auf größte Entfernungen sichtbar zu sein, und dass es sich viel leichter und
besser beobachten lässt als alle Türme und Signale, ja sogar besser als mehrere
zusammengestellte \textsc{Argand}sche Lampen bei Nacht. Um diese Idee praktisch
nutzbar zu machen, brauchte man einen besonderen Apparat oder ein Instrument,
mit dem sich das reflektierte Sonnenlicht ununterbrochen mit größter Genauigkeit
und Sicherheit auf jeden beliebigen noch so weit entfernten Punkt lenken lässt.
Obwohl ich die Vorrichtung für diesen Zweck im Wesentlichen bereits vollständig
entworfen hatte, war es schwierig, sie anfertigen zu lassen. Das galt besonders,
wenn ein auswärtiger Instrumentenbauer die Ausführung hätte übernehmen sollen.
Dann wäre der Vorteil laufender mündlicher Beratung bei einzelnen technischen
Schwierigkeiten entfallen, und die Fertigstellung hätte sich zumindest sehr in
die Länge gezogen. Dieses Problem näherte sich jedoch seiner Lösung, als unser
geschickter Inspektor \textsc{Rumpf}, der während eines großen Teils des Winters
nicht hier gewesen war, gegen Ostern nach Göttingen zurückkehrte und kurz
darauf die Anfertigung eines solchen Instruments \mbox{übernahm.\enspace\mbox{.\ .\ .\ .\ .\ .}}

\par\noindent\hspace*{\parindent}Bei der hannoverschen Gradmessung ist es ein besonders wesentlicher und
wichtiger Umstand, dass sich die ersten von der hiesigen Sternwarte ausgehenden Dreiecksseiten
mit den fest installierten Meridianinstrumenten der Sternwarte selbst aufs
Genaueste orientieren lassen. Der Blick von der Sternwarte in Meridianrichtung
war ursprünglich weder nach Norden noch nach Süden frei. Hätte sich
dieser Mangel nicht beheben lassen, so wäre jener äußerst wichtige Vorteil
gar nicht
\par\clearpage

% Printed 409
\noindent\null\hfill{\small\textbf{ZUR HANNOVERSCHEN TRIANGULATION.}}\hfill\makebox[0pt][r]{\fontsize{13}{16}\selectfont 409}\par\medskip
\noindent vorhanden gewesen, und zugleich wäre dies für immer ein grundlegender
Fehler bei der Wahl des Standorts der Sternwarte geblieben. Glücklicherweise
war die eine Hälfte dieser Schwierigkeit bereits überwunden. Den vorher durch
die Gärten vor Göttingen versperrten Blick nach Norden hatte ich schon im
Herbst 1820 geöffnet und auf einem Berg nahe Weende ein provisorisches
Meridianzeichen errichten lassen. Viel größer waren hingegen die Schwierigkeiten
auf der Südseite des Meridians, wo ein dichter, hoher, drei Stunden entfernter
Wald die Aussicht begrenzte. Dieses Hindernis musste, wenn möglich, überwunden
\mbox{werden.\enspace\mbox{.\ .\ .\ .\ .\ .}}

\par\noindent\hspace*{\parindent}In dieser Zeit hatte Inspektor \textsc{Rumpf} bereits eifrig an dem oben
erwähnten Instrument gearbeitet, das ich im Folgenden mit dem ihm gegebenen
Namen Heliotrop bezeichnen werde. Noch vor dessen Fertigstellung war ich
jedoch auf die Idee gekommen, einen einfachen Spiegelsextanten zu einer
Art Vizeheliotrop umzubauen, der zwar wesentlich weniger vollkommen als
jenes Instrument selbst, bei geschickter Handhabung aber ebenfalls brauchbar
war. Die damit bereits auf Entfernungen von beinahe 2 Meilen durchgeführten
Versuche bestätigten die enorme Stärke des reflektierten Sonnenlichts in einem
Maß, das meine Erwartung fast \mbox{übertraf.\enspace\mbox{.\ .\ .\ .\ .\ .}}
\par\bigskip

\begin{center}
[Aus einem Bericht vom 31.~Januar 1823 »über die Arbeiten der Gradmessung im Jahre 1822«.]
\end{center}

\par\noindent\hspace*{\parindent}1821 war die Triangulation als umfangreichster Teil der gesamten Arbeit
begonnen worden, und die Messungen an fünf Dreieckspunkten waren abgeschlossen,
nämlich in der Göttinger Sternwarte, beim Meridianzeichen, auf dem Hohehagen,
dem Hils und dem Brocken. Für die Fortsetzung waren vier weiter nördlich
gelegene Punkte ausgewählt worden, nämlich Lichtenberg, der Deister, der
Wohlenberg im Amt Gifhorn und ein Berg bei Brelingen in der Amtsvogtei
Bissendorf. Ich selbst hatte diese Punkte jedoch noch nicht besucht, und bei
einigen war noch fraglich, ob sie sich als Dreieckspunkte eignen würden. Dies
musste geklärt werden, bevor sich die ersten Arbeiten des Jahres 1822 planen
ließen.

\par\smallskip
\noindent\hspace*{.030\linewidth}{\footnotesize IX.}\hfill{\small 52}\hspace*{.137\linewidth}\par
\clearpage

% Printed 410
\noindent\makebox[0pt][l]{\fontsize{13}{16}\selectfont 410}\hfill{\small\textbf{NACHLASS.}}\hfill\null\par\medskip
\par\noindent\hspace*{\parindent}Doch das war nur der kleinste Teil der notwendigen Voruntersuchungen.
Die Arbeiten näherten sich nun der Lüneburger Heide, einem ganz flachen
Landstrich, in dem der Mangel an beherrschenden Punkten und die fast
unzähligen größeren und kleineren Waldstücke, die ihn schachbrettartig bedecken,
die Bildung einigermaßen großer Dreiecke außerordentlich erschweren. Ich
kannte diese Schwierigkeiten bereits aus den Berichten des französischen
Obersten \textsc{Epailly}, der 1804 und 1805 die französischen Messungen im
Kurfürstentum Hannover geleitet hatte. Dieser Ingenieur hatte die
Schwierigkeiten des Geländes für so groß gehalten, dass er die Bildung eines
Dreieckssystems von der Aller bis zur Elbe für unmöglich erklärt hatte. Daher
hatte er die Verbindung auf einem gewaltigen Umweg hergestellt, nämlich
durch Dreiecke entlang der Weser bis zu ihrer Mündung und dann wieder
die Elbe hinauf bis Hamburg. Ein solches Verfahren käme allenfalls als
Notbehelf bei einer Landesvermessung infrage, wäre aber bei einer
Gradmessung völlig unzulässig.

\par\noindent\hspace*{\parindent}Diese Umstände machten vor Beginn der eigentlichen Messarbeiten
eine Erkundungsreise notwendig. Diese Reise, zu der ich von meinen drei
Gehilfen nur Hauptmann \textsc{Müller} hinzuzog, beschäftigte mich vom
28.~April bis zum 1.~Juni. Ich übergehe hier die Einzelheiten der mühsamen
Untersuchungen und führe nur ihre wichtigsten Ergebnisse an.

\par\noindent\hspace*{\parindent}Der Brelingerberg erwies sich für die Verbindung als ungeeignet. Das
wurde aber mehr als aufgewogen durch die glückliche Entdeckung, dass sich
zwei neue, noch weiter nördlich gelegene Punkte, ein hoch gelegener Acker
bei Garssen und der Falkenberg, jener fast eine Meile nordöstlich von Celle,
dieser eine Meile nordwestlich von Bergen, beide unmittelbar mit Lichtenberg
und mit dem Deister verbinden ließen. Dadurch wurde der Wohlenberg
überflüssig, und die ersten Messarbeiten standen nun eindeutig und sicher
fest. Was das weitere Vordringen nach Norden betraf, so fand ich die
Schwierigkeiten tatsächlich so groß wie erwartet. Doch es war mir auch
gelungen, gleichsam im Herzen der Heide die Durchführbarkeit zweier guter
Dreiecke zwischen den vier Punkten Falkenberg, Hauselberg (in der Amtsvogtei
Hermannsburg), Wulfsode (im Amt Ebstorf) und Wilsede (an der äußersten
Südwestgrenze des Amts Winsen an der Luhe) festzustellen. Es zeigte
